\documentclass[11pt,a4paper]{article}

\usepackage[T1]{fontenc}
\usepackage[utf8]{inputenc}
\usepackage[margin=2.5cm]{geometry}
\usepackage{amsmath,amssymb,amsfonts,amsthm,mathtools}
\usepackage{bm}
\usepackage{graphicx}
\usepackage{booktabs}
\usepackage{multirow}
\usepackage{xcolor}
\usepackage{enumitem}
\usepackage{hyperref}
\usepackage{url}
\usepackage{xurl}
\usepackage{array}
\usepackage{longtable}
\usepackage{float}
\usepackage{tikz}
\usepackage{pgfplots}
\pgfplotsset{compat=1.18}
\usetikzlibrary{positioning,arrows.meta,shapes.geometric,calc,
                decorations.pathreplacing}
\usepackage{quantikz}
\usepackage{listings}
\usepackage{fancyhdr}
\usepackage{mdframed}
\usepackage{colortbl}
\usepackage{caption}

\providecommand{\dg}{\ensuremath{{}^{\circ}}}
\newcommand{\Ry}{R_y}
\renewcommand{\ket}[1]{|#1\rangle}
\newcommand{\TV}{\operatorname{TV}}
\newcommand{\Fid}{\operatorname{Fid}}
\newcommand{\HS}{H_{\mathrm{S}}}

%% Framed environments
\newmdenv[
  linecolor=black!40,
  backgroundcolor=gray!5,
  roundcorner=4pt,
  innerleftmargin=10pt,
  innerrightmargin=10pt,
  innertopmargin=8pt,
  innerbottommargin=8pt,
]{protocolbox}

\newmdenv[
  linecolor=blue!50!black,
  backgroundcolor=blue!4,
  roundcorner=4pt,
  innerleftmargin=10pt,
  innerrightmargin=10pt,
  innertopmargin=8pt,
  innerbottommargin=8pt,
]{notebox}

\pagestyle{fancy}
\fancyhf{}
\renewcommand{\headrulewidth}{0.4pt}
\fancyhead[L]{\small\textit{Supplementary Material: Extended Experimental Protocol v2.0}}
\fancyhead[R]{\small\textit{Grover--Rudolph Characterisation on SpinQ Triangulum}}
\fancyfoot[C]{\thepage}

\hypersetup{
  colorlinks=true,
  linkcolor=blue!60!black,
  urlcolor=blue!60!black,
  citecolor=blue!60!black
}

\lstset{
  basicstyle=\ttfamily\small,
  breaklines=true,
  frame=single,
  backgroundcolor=\color{gray!8},
  rulecolor=\color{black!30},
  numbers=left,
  numberstyle=\tiny\color{gray},
  numbersep=5pt,
  keywordstyle=\color{blue!70!black}\bfseries,
  commentstyle=\color{green!50!black}\itshape,
  stringstyle=\color{red!60!black},
  language=Python
}

\definecolor{refcolor}{RGB}{220,235,220}
\definecolor{newcolor}{RGB}{210,225,245}

\begin{document}

%%--------------------------------------------------------------
\begin{center}
  {\LARGE\bfseries Supplementary Material}\\[0.6em]
  {\Large Extended Experimental Protocol for the\\[0.2em]
  Grover--Rudolph Hardware Characterisation Campaign\\[0.2em]
  on the SpinQ Triangulum Platform}\\[1.0em]
  {\normalsize Antonio Falcó, Daniela Falcó-Pomares,
   Juan-Carlos Latorre, Yuefeng Feng, \textit{et al.}}\\[0.4em]
  {\small \textit{Companion to:}
  ``Practical Grover--Rudolph state preparation on a 3-qubit NMR
  quantum computer: from rigorous ancilla-free compilation to staged
  experimental diagnostics''}\\[0.3em]
  {\small Version 2.0 --- \today}
\end{center}

\bigskip

%%--------------------------------------------------------------
\begin{abstract}
This document provides the complete experimental protocol for the
extended Grover--Rudolph characterisation campaign on the SpinQ
Triangulum 3-qubit NMR quantum computer.
The campaign extends the first study (Protocol v1.0) in one central
direction: rather than characterising hardware performance for a single
benchmark distribution, it systematically probes the full state-preparation
capability of the device across a heterogeneous suite of seven target
probability distributions covering the full range of Shannon entropy,
symmetry, concentration, and structural complexity that the 3-qubit
Grover--Rudolph construction can address.
This broader characterisation allows for a distribution-dependent
performance map of the Triangulum, separates hardware-generic error
from circuit-structure-specific error, and provides a robust statistical
baseline for cross-platform comparison.
The campaign is designed for automated execution: a parametric pipeline
generates Grover--Rudolph angles, compiles circuits, submits jobs, and
aggregates results for each target distribution without manual
intervention between runs.
The protocol is organised into seven sections:
(\S\ref{sec:prerequisites}) prerequisites and hardware setup;
(\S\ref{sec:distributions}) the target distribution suite;
(\S\ref{sec:gr-engine}) the parametric Grover--Rudolph compilation engine;
(\S\ref{sec:circuits}) circuit structure and ladder variants;
(\S\ref{sec:automation}) automated execution pipeline;
(\S\ref{sec:postprocessing}) data normalization, readout calibration,
and metric computation;
(\S\ref{sec:analysis}) cross-distribution analysis and characterisation outputs.
\end{abstract}

\tableofcontents

\bigskip\hrule\bigskip

%%==============================================================
\section{Prerequisites and hardware setup}
\label{sec:prerequisites}
%%==============================================================

\subsection{Hardware platform}

The experiments are carried out on the \textbf{SpinQ Triangulum},
a commercial 3-qubit desktop NMR quantum computer \cite{FengEtAl2022}.
The device exposes the native gate dictionary
$\{\Ry(\theta), X, \mathrm{CNOT}\}$.
All higher-level operations are decomposed into this set prior to
execution. The campaign has been designed to run on one or two
Triangulum units simultaneously; if two units are available, the
distribution suite can be split across devices to reduce total
wall-clock time (see Section~\ref{sec:automation}).

\begin{protocolbox}
\textbf{Required hardware.}
One or two SpinQ Triangulum units with firmware version recorded
in every run log.\\[4pt]
\textbf{Pre-session check.}
Before each session, verify that a bare measurement (no gates)
returns $\ket{000}$ with probability $\geq 0.98$.
If this threshold is not met, allow $\geq 30$ minutes of thermal
equilibration and repeat.
\end{protocolbox}

\subsection{Software dependencies}

\begin{itemize}
  \item \texttt{numpy} $\geq 1.24$
  \item \texttt{scipy} $\geq 1.10$
  \item \texttt{pandas} $\geq 2.0$
  \item \texttt{jsonlines} $\geq 3.1$
  \item SpinQ SDK (version as shipped with the unit; record in log)
\end{itemize}

\noindent Full implementation:
\begin{center}
  \url{https://github.com/afalco/grover-rudolph-practical-implementation}
\end{center}

\subsection{Qubit register conventions}

The register is $(q_0, q_1, q_2)$ with $q_0$ as the most significant
bit:
\[
  \ket{b_0 b_1 b_2} \leftrightarrow k = 4b_0 + 2b_1 + b_2,
  \quad b_i \in \{0,1\},
\]
giving ordered basis
$\{\ket{000}, \ket{001}, \ket{010}, \ket{011},
   \ket{100}, \ket{101}, \ket{110}, \ket{111}\}$,
indexed $k=0,\ldots,7$.

%%==============================================================
\section{Target distribution suite}
\label{sec:distributions}
%%==============================================================

\subsection{Design rationale}

The first campaign \cite{ProtocolV1} used a single symmetric
benchmark distribution (D1 below). While informative, a single
target cannot distinguish between hardware-generic degradation and
circuit-structure-specific error: a device may perform well on
symmetric, low-contrast targets yet poorly on concentrated or
asymmetric ones. The extended campaign addresses this by probing
seven target distributions that span the relevant axes of variation:

\begin{itemize}
  \item \textbf{Shannon entropy} $\HS(p) = -\sum_k p_k \log_2 p_k \in [0, 3]$:
    from maximally concentrated (near 0) to uniform (3 bits).
  \item \textbf{Symmetry:} symmetric vs.\ monotonic vs.\ bimodal
    vs.\ unstructured.
  \item \textbf{Contrast:} ratio of largest to smallest
    nonzero probability, ranging from 1 (uniform) to $\gg 1$ (concentrated).
  \item \textbf{Support size:} all 8 states active vs.\ effective
    support on a subset.
\end{itemize}

\noindent No two distributions in the suite are close in all three
axes simultaneously, ensuring that the hardware is genuinely
exercised in qualitatively different regimes.

\subsection{The seven target distributions}
\label{subsec:distributions}

Table~\ref{tab:distributions} defines the seven distributions.
Each has been chosen so that its Grover--Rudolph angles are
non-degenerate (no angle collapses to $0\dg{}$ or $180\dg{}$ except
intentionally), and so that the suite as a whole covers the
entropy range $[0.54, 3.00]$ bits with no large gaps.

\begin{table}[H]
\centering
\caption{Target distribution suite. $\HS$: Shannon entropy in bits.
Contrast: $\max_k p_k / \min_{k:\,p_k>0} p_k$.
Distributions D5 and D6 are generated from fixed random seeds
(see \S\ref{subsec:random}) and their values are approximate.}
\label{tab:distributions}
\setlength{\tabcolsep}{4pt}
\begin{tabular}{clp{4.8cm}cc}
\toprule
ID & Type & $p = (p_0,\ldots,p_7)$ & $\HS$ (bits) & Contrast \\
\midrule
\rowcolor{refcolor}
D1 & Symmetric (reference) &
  $(0.05, 0.10, 0.15, 0.20, 0.20, 0.15, 0.10, 0.05)$ &
  2.77 & 4.0 \\
D0 & Uniform &
  $(0.125, 0.125, 0.125, 0.125, 0.125, 0.125, 0.125, 0.125)$ &
  3.00 & 1.0 \\
D2 & Monotone decreasing &
  $(0.35, 0.25, 0.15, 0.10, 0.07, 0.04, 0.03, 0.01)$ &
  2.57 & 35.0 \\
D3 & Bimodal &
  $(0.02, 0.03, 0.05, 0.40, 0.40, 0.05, 0.03, 0.02)$ &
  2.17 & 20.0 \\
D4 & Unimodal / bell &
  $(0.02, 0.05, 0.12, 0.31, 0.31, 0.12, 0.05, 0.02)$ &
  2.53 & 15.5 \\
\rowcolor{newcolor}
D5 & Random (sparse) &
  Dirichlet($\alpha=0.5$), seed 42 &
  $\approx 1.8$ & $\gg 1$ \\
\rowcolor{newcolor}
D6 & Random (generic) &
  Dirichlet($\alpha=1.0$), seed 137 &
  $\approx 2.5$ & $\sim 10$ \\
\bottomrule
\end{tabular}
\end{table}

\noindent\textcolor{blue!60!black}{\textbf{Note.}}
Distributions D5 and D6 (shaded blue) are generated
programmatically from fixed seeds (see \S\ref{subsec:random})
and must be computed once before the campaign begins.
Their exact values must be stored in the campaign configuration
file so that every participant uses identical targets.

\subsection{Rationale for each distribution}

\paragraph{D0 — Uniform.}
The uniform distribution is the trivial fixed point of the
Grover--Rudolph tree: all splitting angles equal $90\dg{}$ and
the UCRy block reduces to a product of single-qubit Hadamard-like
rotations. It provides an absolute baseline and tests whether the
hardware can maintain equal amplitudes in the presence of crosstalk
and relaxation.

\paragraph{D1 — Symmetric (reference).}
The distribution used in the first campaign \cite{ProtocolV1}.
Retained as the inter-campaign anchor point. Any deviation from the
first campaign's results in this distribution can be attributed to
hardware drift rather than experimental design differences.

\paragraph{D2 — Monotone decreasing.}
Strong asymmetry at every level of the dyadic tree: the root angle
($\theta_0 \approx 53\dg{}$) departs significantly from $90\dg{}$,
and the level-1 and level-2 angles are likewise unbalanced. This
stresses the UCRy block with highly non-uniform control rotations
and exercises the hardware in a regime qualitatively different from
D1.

\paragraph{D3 — Bimodal.}
Mass concentrated near two internal states ($\ket{011}$ and
$\ket{100}$). The root split is nearly balanced ($\theta_0 \approx
90\dg{}$), but the level-2 UCRy block must produce very small
rotations ($\theta_{2|00}, \theta_{2|11} \to 0$) alongside large
ones ($\theta_{2|01}, \theta_{2|10} \to 90\dg{}$). This tests
whether the hardware can faithfully execute near-zero controlled
rotations without crosstalk-induced leakage.

\paragraph{D4 — Unimodal / bell.}
A unimodal distribution that is close to D1 in entropy but has a
higher contrast and a more pronounced central peak. It serves as
an intermediate between the symmetric reference and the bimodal
case, and tests the device's ability to suppress tails.

\paragraph{D5 — Random sparse (Dirichlet $\alpha=0.5$).}
A randomly drawn distribution with concentration parameter below 1,
which tends to produce sparse vectors with most mass on a few states.
It has no exploitable symmetry and produces generically unbalanced
trees. The low-$\alpha$ regime is relevant for practical amplitude
encoding scenarios where the target law is irregular.

\paragraph{D6 — Random generic (Dirichlet $\alpha=1.0$).}
A randomly drawn distribution from the uniform prior over the
probability simplex. It serves as a generic test case without
engineered structure. Together with D5, it assesses whether the
performance trends observed in D0--D4 generalize to arbitrary targets.

\subsection{Generating the random distributions}
\label{subsec:random}

Distributions D5 and D6 are generated as follows and must be
computed once, stored, and shared with all participants before the
campaign:

\begin{lstlisting}[caption={Deterministic generation of random distributions D5 and D6.}]
import numpy as np

def dirichlet_dist(alpha: float, seed: int,
                   n: int = 8) -> np.ndarray:
    rng = np.random.default_rng(seed)
    p = rng.dirichlet(np.full(n, alpha))
    # Clip any value below 1e-4 to avoid near-zero entries
    p = np.clip(p, 1e-4, None)
    return p / p.sum()

D5 = dirichlet_dist(alpha=0.5, seed=42)
D6 = dirichlet_dist(alpha=1.0, seed=137)

# Persist to campaign config
import json
config = {"D5": D5.tolist(), "D6": D6.tolist()}
with open("campaign_distributions.json", "w") as f:
    json.dump(config, f, indent=2)
\end{lstlisting}

\begin{protocolbox}
\textbf{Mandatory verification.}
Before submitting any circuit to hardware, all participants must
load \texttt{campaign\_distributions.json} from the shared
repository and verify that their local values of D5 and D6
match to at least 12 decimal places. Use the reference checksums
provided in Appendix~\ref{sec:checksums}.
\end{protocolbox}

%%==============================================================
\section{Parametric Grover--Rudolph compilation engine}
\label{sec:gr-engine}
%%==============================================================

\subsection{General angle formula}

For a general 3-qubit target $p = (p_0, \ldots, p_7)$, the
Grover--Rudolph rotation angles are obtained from the dyadic
conditional masses \cite{GroverRudolph2002,
FalcoFalcoPomaresMatthies2026}. Define the partial sums:
\begin{align}
  s_L &= p_0+p_1+p_2+p_3, \quad s_R = p_4+p_5+p_6+p_7, \\
  s_{LL} &= p_0+p_1, \quad s_{LR} = p_2+p_3, \quad
  s_{RL} = p_4+p_5, \quad s_{RR} = p_6+p_7.
\end{align}
The rotation angles are then:
\begin{align}
  \theta_0 &= 2\arccos\!\sqrt{s_L},
  \label{eq:theta0-gen}\\
  \theta_{1|0} &= 2\arccos\!\sqrt{s_{LL}/s_L}, \qquad
  \theta_{1|1}  = 2\arccos\!\sqrt{s_{RL}/s_R},
  \label{eq:theta1-gen}\\
  \theta_{2|b_0 b_1} &= 2\arccos\!\sqrt{p_{b_0 b_1 0} /
    (p_{b_0 b_1 0} + p_{b_0 b_1 1})},
    \quad (b_0, b_1) \in \{00, 01, 10, 11\},
  \label{eq:theta2-gen}
\end{align}
where $p_{b_0 b_1 b_2}$ denotes $p_k$ with
$k = 4b_0 + 2b_1 + b_2$.
If $p_{b_0 b_1 0} + p_{b_0 b_1 1} = 0$ (zero subtree mass),
set $\theta_{2|b_0 b_1} = 0$ by convention.

\subsection{Walsh--Hadamard transform for the UCRy block}

The four level-2 angles are decomposed into the ancilla-free
ladder representation via the Walsh--Hadamard-type transform
\cite{FalcoFalcoPomaresMatthies2026}. Let
$\alpha_j = \theta_{2|j}/2$ for $j \in \{00, 01, 10, 11\}$
(ordered as binary $j = 0, 1, 2, 3$). The ladder angles are:
\begin{equation}
  \begin{pmatrix} \phi_0 \\ \phi_1 \\ \phi_2 \\ \phi_3 \end{pmatrix}
  = \frac{1}{4}
  \begin{pmatrix}
    1 & 1 & 1 & 1 \\
    1 & -1 & 1 & -1 \\
    1 & 1 & -1 & -1 \\
    1 & -1 & -1 & 1
  \end{pmatrix}
  \begin{pmatrix}
    \alpha_{00} \\ \alpha_{01} \\ \alpha_{10} \\ \alpha_{11}
  \end{pmatrix}.
  \label{eq:wht}
\end{equation}
The full rotation angles commanded to the hardware are
$2\phi_i$ (in degrees).

\subsection{Reference implementation}

\begin{lstlisting}[caption={Parametric GR angle computation for any 3-qubit distribution.}]
import numpy as np

def gr_angles(p: np.ndarray) -> dict:
    """
    Compute all Grover-Rudolph angles for a 3-qubit
    probability vector p of length 8.
    Returns angles in degrees.
    """
    assert len(p) == 8 and np.isclose(p.sum(), 1.0)

    sL  = p[0]+p[1]+p[2]+p[3]
    sR  = 1.0 - sL
    sLL = p[0]+p[1]; sLR = p[2]+p[3]
    sRL = p[4]+p[5]; sRR = p[6]+p[7]

    def safe_arccos(num, den):
        if den < 1e-15:
            return 0.0
        return 2 * np.degrees(np.arccos(np.sqrt(
            np.clip(num/den, 0.0, 1.0))))

    theta0   = safe_arccos(sL, 1.0)
    theta1_0 = safe_arccos(sLL, sL)
    theta1_1 = safe_arccos(sRL, sR)
    theta2   = {
        '00': safe_arccos(p[0], sLL),
        '01': safe_arccos(p[2], sLR),
        '10': safe_arccos(p[4], sRL),
        '11': safe_arccos(p[6], sRR),
    }

    # Walsh-Hadamard transform for UCRy ladder
    alpha = np.array([theta2['00'], theta2['01'],
                      theta2['10'], theta2['11']]) / 2.0
    H = np.array([[1,1,1,1],[1,-1,1,-1],
                  [1,1,-1,-1],[1,-1,-1,1]]) / 4.0
    phi = H @ alpha  # ladder half-angles in degrees

    return {
        'theta0':   theta0,
        'theta1_0': theta1_0,
        'theta1_1': theta1_1,
        'theta2':   theta2,
        'phi':      phi,          # ladder half-angles
        'ladder_angles': 2*phi,   # full angles for hardware
    }
\end{lstlisting}

\begin{protocolbox}
\textbf{Verification step (mandatory for every distribution).}
After computing angles, run the ideal simulator and confirm
$\TV(q_{\mathrm{sim}}, p) \leq 10^{-10}$ and
$\Fid(q_{\mathrm{sim}}, p) = 1.000$.
Any failure indicates a formula or implementation error;
do not proceed to hardware until resolved.
\end{protocolbox}

%%==============================================================
\section{Circuit structure and ladder variants}
\label{sec:circuits}
%%==============================================================

\subsection{Diagnostic stages}

The three nested diagnostic stages $L0$, $L01$, FULL are preserved
from the first campaign and applied identically to every
distribution in the suite.
Their definitions in terms of the dyadic tree levels are unchanged:
$L0$ applies only $\theta_0$; $L01$ applies $\theta_0$ and the
level-1 angles; FULL applies the complete construction.

The gate counts depend only on the stage and ladder, not on the
specific distribution, since the circuit topology is fixed by the
UCRy decomposition:

\begin{table}[H]
\centering
\caption{Gate counts per stage and ladder (identical for all
  distributions D0--D6).}
\label{tab:gatecounts}
\begin{tabular}{lcccc}
\toprule
Stage & Total gates & $\Ry$ & CNOT & $X$ \\
\midrule
$L0$    &  3 &  3 & 0 & 0 \\
$L01$   & 13 &  7 & 4 & 2 \\
FULL-A  & 21 & 11 & 8 & 2 \\
FULL-B  & 21 & 11 & 8 & 2 \\
\bottomrule
\end{tabular}
\end{table}

\subsection{Ladder variants A and B}

The two FULL circuits differ only in the Gray-code ordering of the
two-control UCRy block on $q_2$, as in the first campaign:

\begin{itemize}
  \item \textbf{Ladder A.} UCRy control sequence:
    $(q_1, q_0, q_0, q_0)$.
    Sub-rotation order:
    $\Ry(2\phi_0)$, CNOT($q_1$),
    $\Ry(2\phi_3)$, CNOT($q_0$),
    $\Ry(2\phi_2)$, CNOT($q_0$),
    $\Ry(2\phi_1)$, CNOT($q_0$).

  \item \textbf{Ladder B.} UCRy control sequence:
    $(q_1, q_1, q_0, q_0)$.
    Sub-rotation order:
    $\Ry(2\phi_0)$, CNOT($q_1$),
    $\Ry(2\phi_1)$, CNOT($q_1$),
    $\Ry(2\phi_2)$, CNOT($q_0$),
    $\Ry(2\phi_3)$, CNOT($q_0$).
\end{itemize}

Note that the ladder angles $\phi_i$ now depend on the target
distribution through equation~\eqref{eq:wht}, so ladders A and B
produce physically different pulse schedules for each distribution
in the suite. The unitary equivalence is preserved in all cases.

%%==============================================================
\section{Automated execution pipeline}
\label{sec:automation}
%%==============================================================

\subsection{Campaign structure}

The full campaign is organized as follows:

\begin{itemize}
  \item \textbf{Distributions:} 7 (D0--D6).
  \item \textbf{Stages per distribution:} 3 ($L0$, $L01$, FULL).
  \item \textbf{Ladder variants (FULL only):} 2 (A, B).
  \item \textbf{Runs per (distribution, stage, ladder):} 25.
  \item \textbf{Total runs:} $7 \times (50 + 50 + 25 + 25) = 1050$.
\end{itemize}

Compared with the first campaign (150 runs, 1 distribution), the
extended campaign is a factor-7 scale-up in total hardware time,
motivating the automated pipeline described below.

\subsection{Pipeline architecture}

The pipeline is organised as four independent modules that are
called sequentially for each (distribution, stage, ladder) triplet:

\begin{enumerate}[label=\arabic*.]
  \item \textbf{Angle engine.} Loads the target distribution,
    calls \texttt{gr\_angles()}, and returns the full angle dictionary.
  \item \textbf{Circuit compiler.} Takes the angle dictionary and
    stage/ladder specification, instantiates the gate sequence, and
    serializes it to the platform's JSON circuit format.
  \item \textbf{Job executor.} Submits the compiled circuit to the
    Triangulum SDK, polls until completion, retrieves raw counts,
    and appends the result to the JSONL log.
  \item \textbf{Run validator.} Checks that the returned count sum
    is within $1\%$ of the nominal shot number and that no
    error status was returned; marks the run as \texttt{failed}
    if either condition is violated.
\end{enumerate}

\subsection{Orchestrator}

\begin{lstlisting}[caption={Campaign orchestrator: iterates over all
  distributions, stages, and ladders.}]
import json
import itertools

DISTRIBUTIONS = {
    'D0': np.full(8, 0.125),
    'D1': np.array([0.05,0.10,0.15,0.20,
                    0.20,0.15,0.10,0.05]),
    'D2': np.array([0.35,0.25,0.15,0.10,
                    0.07,0.04,0.03,0.01]),
    'D3': np.array([0.02,0.03,0.05,0.40,
                    0.40,0.05,0.03,0.02]),
    'D4': np.array([0.02,0.05,0.12,0.31,
                    0.31,0.12,0.05,0.02]),
    'D5': load_from_config('D5'),   # from campaign_distributions.json
    'D6': load_from_config('D6'),
}
STAGES   = ['L0', 'L01', 'FULL']
LADDERS  = {'L0': ['A'], 'L01': ['A'], 'FULL': ['A', 'B']}
N_RUNS   = 25   # runs per (dist, stage, ladder)
LOGFILE  = 'campaign_v2_runs.jsonl'

def run_campaign():
    run_index = 0
    for dist_id, p in DISTRIBUTIONS.items():
        angles = gr_angles(p)
        verify_simulation(p, angles)       # mandatory check
        for stage in STAGES:
            for ladder in LADDERS[stage]:
                circuit = compile_circuit(angles, stage, ladder)
                for rep in range(1, N_RUNS + 1):
                    run_index += 1
                    record = execute_and_log(
                        circuit=circuit,
                        dist_id=dist_id,
                        p_target=p,
                        stage=stage,
                        ladder=ladder,
                        repeat=rep,
                        run_index=run_index,
                    )
                    log_run(record, LOGFILE)

if __name__ == '__main__':
    run_campaign()
\end{lstlisting}

\subsection{Execution order and drift mitigation}

To mitigate systematic hardware drift across the campaign, runs
are interleaved across distributions rather than executing all
runs for one distribution before moving to the next. The
recommended block structure is:

\begin{enumerate}[label=\arabic*.]
  \item For each stage ($L0$, $L01$, FULL), execute one block of 5 runs
    per distribution per ladder, cycling through all 7 distributions.
  \item Repeat the full cycle 5 times to accumulate 25 runs per group.
  \item After every 100 runs, perform a bare-state fidelity check;
    suspend if $P(\ket{000}) < 0.97$.
\end{enumerate}

This interleaved structure ensures that any slow drift in the
device affects all distributions approximately equally, so
inter-distribution comparisons remain valid.

\subsection{Run log schema}

Each run is logged as a single JSON object appended to
\texttt{campaign\_v2\_runs.jsonl}. In addition to the fields
from Protocol~v1.0, the following fields are added:

\begin{longtable}{lll}
\toprule
Field & Type & Description \\
\midrule
\texttt{dist\_id}     & string  & Distribution identifier (D0--D6) \\
\texttt{dist\_type}   & string  & Human-readable type (e.g.\ \texttt{uniform}) \\
\texttt{p\_target}    & array   & Target probability vector (8 floats) \\
\texttt{shannon\_entropy} & float & $\HS(p)$ in bits \\
\texttt{contrast}     & float   & $\max_k p_k / \min_{k:p_k>0} p_k$ \\
\texttt{angles}       & object  & Full angle dictionary from \texttt{gr\_angles()} \\
\texttt{run\_index}   & integer & Global run index across the full campaign \\
\bottomrule
\end{longtable}

%%==============================================================
\section{Data normalization, readout calibration, and metrics}
\label{sec:postprocessing}
%%==============================================================

\subsection{Probability normalization and marginals}

Identical to Protocol v1.0. Raw counts are normalized to
probability vectors via equation~\eqref{eq:normalize}; per-qubit
marginals are computed via equation~\eqref{eq:marginals}.

\begin{equation}
  q_k = \frac{c_k}{\sum_{j=0}^{7} c_j}, \quad k = 0, \ldots, 7.
  \label{eq:normalize}
\end{equation}
\begin{equation}
  P(q_j = b) = \sum_{\substack{k=0\\ b_j(k)=b}}^{7} q_k.
  \label{eq:marginals}
\end{equation}

\subsection{Readout calibration}

The $L0$ stage provides a per-qubit assignment error estimate for
each distribution run. Since the $L0$ state is identical for all
distributions (only $q_0$ is active), the readout error
estimates from D0--D6 can be pooled to obtain a more precise
characterisation of the device's measurement bias across the
campaign:
\begin{equation}
  \varepsilon_{0\to1}(Q_j) =
    1 - \frac{1}{7 \cdot 50}
    \sum_{\text{dist}} \sum_{\text{run}} P_{\mathrm{exp}}(q_j=0)
    \Big|_{L0, \text{dist}},
  \quad j \in \{1, 2\}.
  \label{eq:readout-pooled}
\end{equation}
This pooled estimate uses $7 \times 50 = 350$ $L0$ runs and
gives a substantially tighter characterisation of readout error
than the 50 runs available in the first campaign.

\subsection{Run-level metrics}

For each run, compute the six metrics as in Protocol v1.0:

\begin{align}
  \TV(p, q) &= \tfrac{1}{2} \sum_{k=0}^{7} |p_k - q_k|,
  \label{eq:tv}\\
  \|p-q\|_2 &= \bigl(\textstyle\sum_{k=0}^{7}(p_k-q_k)^2\bigr)^{1/2},
  \label{eq:l2}\\
  \Fid(p, q) &= \bigl(\textstyle\sum_{k=0}^{7}\sqrt{p_k q_k}\bigr)^2.
  \label{eq:fidelity}
\end{align}

Compute against both the simulator output and the target:
\texttt{tv\_vs\_sim}, \texttt{l2\_vs\_sim}, \texttt{fidelity\_vs\_sim},
\texttt{tv\_vs\_target}, \texttt{l2\_vs\_target},
\texttt{fidelity\_vs\_target}.

\subsection{Output data artifacts}

The campaign produces the following artifacts:

\begin{enumerate}[label=\arabic*.]
  \item \textbf{\texttt{campaign\_v2\_runs.jsonl}} --- complete
    per-run records for all 1050 runs.
  \item \textbf{\texttt{runs\_flat\_v2.csv}} --- one row per run,
    all metric and probability fields, with additional columns
    \texttt{dist\_id}, \texttt{shannon\_entropy}, \texttt{contrast}.
  \item \textbf{\texttt{summary\_by\_dist\_stage.csv}} --- one row
    per (dist\_id, stage, ladder) group ($7 \times 4 = 28$ rows),
    with means and standard deviations of all metrics.
  \item \textbf{\texttt{summary\_by\_dist\_FULL.csv}} --- one row
    per distribution at the FULL stage only ($7 \times 2 = 14$ rows),
    used for the cross-distribution characterisation analysis.
  \item \textbf{\texttt{circuits\_flat\_v2.csv}} --- one row per
    (dist\_id, stage, ladder) combination (28 rows), with angle
    values, gate counts, and circuit payloads.
\end{enumerate}

%%==============================================================
\section{Cross-distribution analysis and characterisation outputs}
\label{sec:analysis}
%%==============================================================

\subsection{Primary characterisation plot}

The central output of the extended campaign is a
\emph{performance map} of the Triangulum: a plot of mean FULL-stage
classical fidelity vs.\ Shannon entropy $\HS(p)$, with one point per
distribution and error bars equal to the standard deviation over 25
runs per ladder. Ladders A and B are shown separately.
This plot provides a quantitative answer to the question of how
hardware performance varies with target complexity.

\begin{protocolbox}
\textbf{Expected qualitative behaviour.}
Based on the first campaign and the structure of the circuits,
we anticipate a monotone-increasing trend (higher entropy
$\Rightarrow$ less demanding UCRy block $\Rightarrow$ higher
fidelity), with D0 (uniform) achieving the highest fidelity and
D3 or D5 (bimodal / sparse) the lowest. Deviations from this
trend are scientifically informative and should be reported and
discussed.
\end{protocolbox}

\subsection{Stage-wise degradation per distribution}

For each distribution, compute the stage-wise TV distances and
fidelities as in the first campaign, and report the degradation
trajectory $(L0 \to L01 \to \mathrm{FULL})$.
A distribution for which the $L0 \to L01$ step is larger than the
$L01 \to \mathrm{FULL}$ step (in TV units) exhibits primarily
level-1 control error; the reverse pattern implicates the UCRy
block.

\subsection{Ladder A vs.\ B comparison per distribution}

Apply the two-sided Mann--Whitney $U$ test for ladders A vs.\ B
at the FULL stage, independently for each distribution and each
metric. Summarise results in a $7 \times 3$ table (distributions
$\times$ metrics) of $p$-values. Test for a consistent pattern:
if ladder ordering matters more for concentrated distributions
(D3, D5) than for smooth ones (D0, D1), this suggests that
near-zero controlled rotations are particularly sensitive to
Gray-code ordering.

\subsection{Readout error characterisation across distributions}

Plot $\Delta\TV = \TV(q_{\mathrm{exp}}, p) -
\TV(q_{\mathrm{exp}}, q_{\mathrm{sim}})$ at stage $L0$ as a
function of distribution. Since the $L0$ circuit is identical
for all distributions, any variation in $\Delta\TV$ across
distributions is attributable to run-to-run hardware variability,
not to the target itself. This provides a direct estimate of
the precision of the readout error characterisation.

\subsection{Regression: fidelity vs.\ circuit properties}

Fit a simple linear model:
\begin{equation}
  \Fid_{\mathrm{FULL}} \approx
    \beta_0
    + \beta_1 \HS(p)
    + \beta_2 \log\mathrm{Contrast}(p)
    + \beta_3 \max_j |\theta_{2|j} - 90\dg{}|
    + \varepsilon.
  \label{eq:regression}
\end{equation}
Report the estimated coefficients, $R^2$, and the 95\% confidence
intervals. This model, while approximate, quantifies the relative
contribution of entropy, contrast, and UCRy angle extremity to
hardware-level fidelity loss.

\begin{lstlisting}[caption={Cross-distribution regression.}]
import numpy as np
from scipy import stats

def characterisation_regression(summary_df):
    """
    summary_df: one row per distribution at FULL stage.
    Columns: shannon_entropy, log_contrast,
             max_ucry_deviation, fidelity_vs_target_mean
    """
    X = summary_df[['shannon_entropy',
                     'log_contrast',
                     'max_ucry_deviation']].values
    X = np.column_stack([np.ones(len(X)), X])
    y = summary_df['fidelity_vs_target_mean'].values

    beta, res, rank, sv = np.linalg.lstsq(X, y,
                                            rcond=None)
    y_hat = X @ beta
    ss_res = np.sum((y - y_hat)**2)
    ss_tot = np.sum((y - y.mean())**2)
    r2 = 1 - ss_res / ss_tot

    return {'beta': beta, 'R2': r2}
\end{lstlisting}

%%==============================================================
\section*{Data and code availability}
%%==============================================================

All data artifacts and the full automation pipeline will be
released at:
\begin{center}
  \url{https://github.com/afalco/grover-rudolph-practical-implementation}
\end{center}
The repository will include a \texttt{campaigns/v2/} subdirectory
with the orchestrator, angle engine, and post-processing scripts,
as well as the fixed campaign configuration file
\texttt{campaign\_distributions.json}.

%%==============================================================
\appendix

\section{Distribution checksums}
\label{sec:checksums}

After generating D5 and D6 via the procedure in
Section~\ref{subsec:random}, verify the following SHA-256 checksums
of the JSON-serialised vectors (compact format, 15 decimal places):

\medskip
\noindent\texttt{D5:} \texttt{[to be completed after generation]}\\
\texttt{D6:} \texttt{[to be completed after generation]}

\medskip
\noindent These values will be committed to the repository before
the campaign begins and serve as the reference for all participants.

\section{Expected results: first campaign anchor (D1)}
\label{sec:anchor}

For distribution D1 (the reference from Protocol v1.0), the
extended campaign should reproduce the following results within
$\pm 2\sigma$ of the first campaign values. Systematic deviations
indicate hardware drift between campaigns.

\begin{table}[H]
\centering
\caption{First-campaign reference values for D1 (25 runs per ladder per stage).}
\begin{tabular}{lcccc}
\toprule
Stage & Ladder & TV vs target & $\ell^2$ vs target & Fidelity vs target \\
\midrule
\multirow{2}{*}{$L0$}
  & A & $0.522\pm0.039$ & $0.455\pm0.031$ & $0.657\pm0.043$ \\
  & B & $0.530\pm0.041$ & $0.465\pm0.032$ & $0.647\pm0.043$ \\
\multirow{2}{*}{$L01$}
  & A & $0.352\pm0.038$ & $0.299\pm0.030$ & $0.804\pm0.036$ \\
  & B & $0.345\pm0.046$ & $0.290\pm0.035$ & $0.817\pm0.048$ \\
\multirow{2}{*}{FULL}
  & A & $0.234\pm0.047$ & $0.195\pm0.037$ & $0.916\pm0.033$ \\
  & B & $0.232\pm0.059$ & $0.201\pm0.047$ & $0.909\pm0.038$ \\
\bottomrule
\end{tabular}
\end{table}

%%==============================================================
\begin{thebibliography}{9}

\bibitem{GroverRudolph2002}
Grover, L., Rudolph, T.: Creating superpositions that correspond to
efficiently integrable probability distributions.
arXiv:quant-ph/0208112 (2002)

\bibitem{FalcoFalcoPomaresMatthies2026}
Falcó, A., Falcó-Pomares, D., Matthies, H.G.:
A rigorous and self-contained proof of the Grover--Rudolph state
preparation algorithm. arXiv:2601.17930 (2026).
Submitted to AIMS Mathematics

\bibitem{FengEtAl2022}
Feng, G., Hou, S.-Y., Zou, H., et al.:
SpinQ Triangulum: a commercial three-qubit desktop quantum computer.
arXiv:2202.02983 (2022)

\bibitem{ProtocolV1}
Falcó, A., Falcó-Pomares, D., Latorre, J.-C., et al.:
Supplementary Material: Experimental Protocol for the First
Grover--Rudolph Campaign on the SpinQ Triangulum Platform.
Version 1.0 (2026). Available at the companion repository.

\end{thebibliography}

\end{document}
